
\def \Subject {تشخیص احساسات از گفتار}
\def \Session { یک}
% \setcounter{chapter}{\Session-1}
\renewcommand{\chaptername}{سوال}
\chapter{\Subject}
\section{مجموعه داده}
مجموعه داده از سایت 
\lr{kaggle}
دانلود شد و بارگذاری شد. برای دیتالودر تنها اطلاعات گوینده و برچسب هر داده ذخیره شد زیرا خود جمله در اینجا اهمیتی ندارد. اهمیت گوینده نیز در فیلتر کردن مجموعه داده مطابق خواست سوال است. پس از فیلتر کردن داده‌ها کلاس 
\lr{Dataset}
هنگامی که یک داده درخواست می‌شود، فایل مربوط به داده را می‌خواند، طول آن را مطابق کانفیگ تغییر می‌دهد و ویژگی‌های آن را استخراج می‌کند و  سپس ویژگی‌‌ها و برچسب را خروجی می‌دهد. 

بررسی شد که در بین داده‌ها پس از فیلتر کردن، ناهماهنگی خاصی در بین احساسات نباشد. 

\begin{figure}[H]
	\centering
	\includegraphics[width=0.8\textwidth]{images/emotion_distribution.png}
	\caption{توزیع احساسات در بین مجموعه داده‌های مختلف}
	\label{fig:data_dist}
\end{figure}

\begin{figure}[H]
	\centering
	\includegraphics[width=0.8\textwidth]{images/actor_distribution.png}
	\caption{افراد موجود در هر بخش داده}
	\label{fig:actor_data_dist}
\end{figure}

\section{ویژگی‌ها}
برای استخراج ویژگی 
\lr{Mel Spectrogram}
از تابع آماده موجود در
\lr{torch}
و برای بارگذاری مدل 
\lr{Hubert}
از کتابخانه 
\lr{transformers}
استفاده شد. 

برای 
\lr{Mel}
پس از استخراج ویژگی‌های 
\lr{Mel}
باید از آن ها لگاریتم بگیریم تا برای انسان معنی دار تر باشند. برای این کار از تبدیل 
\lr{Amplitude to decibel}
استفاده شد. 

مدل 
\lr{hubert}
نیز پیچیدگی خاصی ندارد. هر ورودی را ابتدا به 
\lr{Wav2Vec2}
می‌دهیم و سپس خروجی آن را به مدل 
\lr{Hubert}
می‌دهیم
و خروجی لایه پنهان اخر را به عنوان نمایش صوت خروجی می‌دهیم. 


\section{طبقه‌بندی}
مشابه خواسته سوال از یک مدل پیچشی برای 
\lr{mel}
و از یک مدل 
\lr{MLP}
برای 
\lr{Hubert}
استفاده می‌کنیم. تعداد کانال‌های مدل پیچشی در جدول 
\ref{tab:conv}
آمده است .

\begin{table}[h!]
	\centering
	\caption{تعداد کانال‌های هر لایه مدل پیچشی}
	\label{tab:conv}
	\begin{tabular}{|r|c|}
		\hline
		شماره لایه & تعداد کانال \\
		\hline
	۱ & ۱۶ \\
	۲ & ۳۲ \\
	۳ & ۶۴ \\
		\hline
	\end{tabular}
	\caption{تعداد کانال‌های هر لایه پیچشی}
\end{table}

برای مدل تماما متصل نیز از یک مدل با یک لایه پنهان استفاده شد. 

\section{آموزش و ارزیابی مدل}
آموزش مدل با همان پارامترهای پیشنهادی انجام شد. تنها برای تعداد دور از مقدار ۲۰ استفاده شد. زیرا با استفاده از ۵ دور، به طور واضحی مدل هنوز اشباع نشده بود. 
\begin{figure}[H]
	\centering
	\includegraphics[width=0.8\textwidth]{images/training_curves_hubert.png}
	\caption{نمودار خطا و دقت مدل روی داده‌های آموزش و صحت‌سنجی با \lr{hubert}}
	\label{fig:loss_acc_plot}
\end{figure}

\begin{figure}[H]
	\centering
	\includegraphics[width=0.8\textwidth]{images/training_curves_mel_spectrogram.png}
	\caption{نمودار خطا و دقت مدل روی داده‌های آموزش و صحت‌سنجی با \lr{mel spectro}}
	\label{fig:loss_acc_plot_mel}
\end{figure}

\begin{figure}[H]
	\centering
	\includegraphics[width=0.8\textwidth]{images/confusion_matrix.png}
	\caption{ماتریس درهم ریختگی}
	\label{fig:confusion}
\end{figure}

روی بهترین مدل توانست به دقت ۷۷ درصد روی داده‌های آزمون برسد. 

با مقایسه نمودارهای 
\ref{fig:loss_acc_plot}
و
\ref{fig:loss_acc_plot_mel}
می‌توان دید که مدل 
\lr{hubert}
هم در عدم بیش‌برازش و هم در دقت به مراتب بهتر از مدل 
\lr{mel spectorgram}
بوده است. مدل 
\lr{mel}
هم اختلاف زیادی بین دقت روی آزمون و آزمایش دارد و هم نتوانسته در دقت به مدل 
\lr{hubert}
برسد. همچنین بعد از دورهای کمی اشباع می‌شود در حالی که مدل 
\lr{hubert}
حتی پس از ۲۰ دور نیز اشباع نشده است. 

\section{سوال امتیازی}
\subsection{معماری مدل \lr{Speech T5}}
مدل در اینجا شامل یک مدل کدگذار کدگشای اصلی همراه با شش بلوک مجزا است. این بلوک‌های مجزا پیچیدگی‌های 
\lr{modality}
ها را هندل می‌کنند و شبکه کدگذار کدگشا نیز اصل پردازش را انجام می‌دهد. 


\begin{figure}[H]
	\centering
	\includegraphics[width=0.8\textwidth]{images/architecture.png}
	\caption{معماری مدل که در مقاله آمده است}
	\label{fig:speech_arg}
\end{figure}

همانطور که درشکل 
\ref{fig:speech_arg}
می‌توان دید،  ابتدا متن و گفتار هر کدام به یک شبکه مجزا داده می‌شوند سپس خروجی این دو به مدل کدگذار کدگشا داده می‌شود. هنگام انتقال اطلاعات از کدگذار به کدگشا، خروجی دو بلوک دیگر که یکی برای متن و یکی برای گفتار است نیز به این اطلاعات اضافه می‌شود. پس از پردازش اطلاعات در کدگشا، خروجی به یک بلوک پس پردازش متن و یک بلوک پس پردازش گفتار داده می‌شود. در ادامه وظیفه هرکدام از این بلوک‌ها توضیح داده شده است. 

\begin{itemize}
	\item \lr{\textbf{Speech Encoder Pre-Net} } 
	
		تبدیل صوت به بردارهای 
		\lr{representation}
		با استفاده از ۷ بلوک 
		\lr{convolutional}
		 برای ورود به رمزگذار.
	
	\item \lr{\textbf{Text Encoder Pre-Net} } 
	
	تبدیل  توکن‌های متن به بردارهای 
	\lr{embedding}
	برای ورود به رمزگذار.
	
	\item \lr{\textbf{Speech Decoder Pre-Net} } 
	
	تبدیل صوت
	 به فضای یکپارچه برای ورود به رمزگشا.
	
	\item \lr{\textbf{Text Decoder Pre-Net}} 
	
	تبدیل توکن‌های متن هدف  به بردارهای 
	\lr{embedding}
	فضای مشترک برای ورود به رمزگشا.
	
	\item \lr{\textbf{Speech Decoder Post-Net} } 
	
	تبدیل خروجی رمزگشا به صوت
برای تولید گفتار 
	
	\item \lr{\textbf{Text Decoder Post-Net}} 
	
	تبدیل خروجی رمزگشا به توزیع احتمالی روی توکن‌های متن با استفاده از یک لایه خطی به همراه تابع 
	\lr{Softmax}
	 برای تولید متن.
\end{itemize}















\subsection{\lr{Pretext task} داده‌های گفتاری مدل}

مدل SpeechT5 با استفاده از دو 
\lr{pretext task}
بر روی داده‌های گفتاری بدون برچسب آموزش داده می‌شود. شیوه اموزش مدل تقریبا شبیه به \lr{bert} است یعنی بخش‌هایی از ورودی ماسک می‌شود و مدل باید آن‌ها را حدس بزند. 

در تسک اول مدل باید صرفا واحد صوتی ماسک شده را پیدا کند. یعنی به نوعی شبیه به پیشبینی توکن در فضای متنی. در اینجا صرفا بخش کدگذار و سمت ورودی مدل آموزش داده می‌شود. 

در تسک دوم مدل باید سعی کند صوت را بازسازی کند یعنی هم بخش ورودی و کدگذار آموزش داده می‌شود و هم بخش خروجی و کدگشا برای تولید صوت.

\subsection{  تسک‌های انجام‌شده با \lr{SpeechT5}}

مدل 
\lr{SpeechT5}
 بر روی شش تسک مختلف  ارزیابی شده است.

\begin{table}[h!]
	\centering
	\begin{tabular}{|p{2.5cm}|p{3.5cm}|p{3.5cm}|}
		\hline
		\textbf{تسک} & \textbf{ورودی} & \textbf{خروجی} \\
		\hline
		تشخیص خودکار گفتار &صوت گفتار & دنبالهٔ کاراکترهای متن \\
		\hline
		تبدیل متن به گفتار & متن& ویژگی‌های Log Mel-Filterbank \\
		\hline
		ترجمه گفتار & صوت به زبان انگلیسی &متن به زبان آلمانی یا فرانسوی \\
		\hline
		تبدیل صدا  & صوت گوینده مبدا & صوت گوینده هدف \\
		\hline
		بهبود گفتار & صوت نویزی & صوت تمیز \\
		\hline
		تشخیص گوینده  & صوت یک گوینده & برچسب گوینده \\
		\hline
	\end{tabular}
	\caption{خلاصه تسک‌های انجام‌شده با SpeechT5 و ساختار ورودی/خروجی آن‌ها}
\end{table}



